\section{Integration by parts}

Integration by parts is the method that reverses the product rule. It is useful when an integrand is a product of two types of functions and substitution does not directly simplify the expression.

The product rule says
\[
(uv)'=u'v+uv'.
\]
Rearranging and integrating gives the integration by parts formula.

\begin{theorembox}
If \(u=u(x)\) and \(v=v(x)\) are differentiable, then
\[
\int u\,dv=uv-\int v\,du.
\]
Equivalently,
\[
\int u(x)v'(x)\,dx=u(x)v(x)-\int v(x)u'(x)\,dx.
\]
\end{theorembox}

The method replaces one integral by another. It is useful only if the new integral is simpler.

\subsection*{A first example}

\begin{examplebox}
Compute
\[
\int x e^x\,dx.
\]
Choose
\[
u=x,
\qquad
dv=e^x\,dx.
\]
Then
\[
du=dx,
\qquad
v=e^x.
\]
By integration by parts,
\[
\int x e^x\,dx=x e^x-\int e^x\,dx.
\]
Therefore
\[
\int x e^x\,dx=x e^x-e^x+C.
\]
\end{examplebox}

The choice of \(u\) matters. We chose \(u=x\) because differentiating it makes it simpler.

\subsection*{Logarithmic functions}

Integration by parts is often useful when a logarithm appears.

\begin{examplebox}
Compute
\[
\int \ln x\,dx,
\qquad x>0.
\]
Although this does not visibly look like a product, we can write
\[
\ln x=1\cdot\ln x.
\]
Choose
\[
u=\ln x,
\qquad
dv=dx.
\]
Then
\[
du=\frac1x\,dx,
\qquad
v=x.
\]
Thus
\[
\int \ln x\,dx=x\ln x-\int x\cdot\frac1x\,dx.
\]
So
\[
\int \ln x\,dx=x\ln x-\int 1\,dx=x\ln x-x+C.
\]
\end{examplebox}

This example shows that a factor of \(1\) may be introduced when it helps us apply the method.

\subsection*{Repeated integration by parts}

Sometimes integration by parts must be applied more than once.

\begin{examplebox}
Compute
\[
\int x^2 e^x\,dx.
\]
Choose
\[
u=x^2,
\qquad
dv=e^x\,dx.
\]
Then
\[
du=2x\,dx,
\qquad
v=e^x.
\]
Therefore
\[
\int x^2e^x\,dx=x^2e^x-\int 2xe^x\,dx.
\]
Now apply integration by parts to \(\int xe^x\,dx\), which was computed earlier:
\[
\int xe^x\,dx=xe^x-e^x+C.
\]
Thus
\[
\int x^2e^x\,dx=x^2e^x-2(xe^x-e^x)+C.
\]
So
\[
\int x^2e^x\,dx=e^x(x^2-2x+2)+C.
\]
\end{examplebox}

The pattern is clear: differentiating the polynomial lowers its degree until the remaining integral is elementary.

\subsection*{Choosing \(u\)}

There is no perfect rule for choosing \(u\), but a useful principle is to choose \(u\) so that differentiating it simplifies the expression. Logarithms, inverse trigonometric functions, and polynomial factors are often good candidates for \(u\).

\begin{summarybox}
When using integration by parts, ask:
\begin{itemize}
\item Does the integrand look like a product?
\item Which factor becomes simpler when differentiated?
\item Which factor can be integrated to find \(v\)?
\item Is the new integral simpler than the original one?
\item Do I need to apply the method repeatedly?
\item Can I check the final answer by differentiating?
\end{itemize}
\end{summarybox}

Integration by parts is not a formula to apply mechanically. It is a way to move differentiation from one factor to another.